% Options for packages loaded elsewhere
\PassOptionsToPackage{unicode}{hyperref}
\PassOptionsToPackage{hyphens}{url}
%
\documentclass[
  ignorenonframetext,
  aspectratio=169,
]{beamer}
\usepackage{pgfpages}
\setbeamertemplate{caption}[numbered]
\setbeamertemplate{caption label separator}{: }
\setbeamercolor{caption name}{fg=normal text.fg}
\beamertemplatenavigationsymbolshorizontal
% Prevent slide breaks in the middle of a paragraph
\widowpenalties 1 10000
\raggedbottom
\setbeamertemplate{part page}{
  \centering
  \begin{beamercolorbox}[sep=16pt,center]{part title}
    \usebeamerfont{part title}\insertpart\par
  \end{beamercolorbox}
}
\setbeamertemplate{section page}{
  \centering
  \begin{beamercolorbox}[sep=12pt,center]{part title}
    \usebeamerfont{section title}\insertsection\par
  \end{beamercolorbox}
}
\setbeamertemplate{subsection page}{
  \centering
  \begin{beamercolorbox}[sep=8pt,center]{part title}
    \usebeamerfont{subsection title}\insertsubsection\par
  \end{beamercolorbox}
}
\AtBeginPart{
  \frame{\partpage}
}
\AtBeginSection{
  \ifbibliography
  \else
    \frame{\sectionpage}
  \fi
}
\AtBeginSubsection{
  \frame{\subsectionpage}
}

\usepackage{amsmath,amssymb}
\usepackage{iftex}
\ifPDFTeX
  \usepackage[T1]{fontenc}
  \usepackage[utf8]{inputenc}
  \usepackage{textcomp} % provide euro and other symbols
\else % if luatex or xetex
  \usepackage{unicode-math}
  \defaultfontfeatures{Scale=MatchLowercase}
  \defaultfontfeatures[\rmfamily]{Ligatures=TeX,Scale=1}
\fi
\usepackage{lmodern}
\ifPDFTeX\else  
    % xetex/luatex font selection
\fi
% Use upquote if available, for straight quotes in verbatim environments
\IfFileExists{upquote.sty}{\usepackage{upquote}}{}
\IfFileExists{microtype.sty}{% use microtype if available
  \usepackage[]{microtype}
  \UseMicrotypeSet[protrusion]{basicmath} % disable protrusion for tt fonts
}{}
\usepackage{xcolor}
\newif\ifbibliography
\setlength{\emergencystretch}{3em} % prevent overfull lines
\setcounter{secnumdepth}{-\maxdimen} % remove section numbering


\providecommand{\tightlist}{%
  \setlength{\itemsep}{0pt}\setlength{\parskip}{0pt}}\usepackage{longtable,booktabs,array}
\usepackage{calc} % for calculating minipage widths
\usepackage{caption}
% Make caption package work with longtable
\makeatletter
\def\fnum@table{\tablename~\thetable}
\makeatother
\usepackage{graphicx}
\makeatletter
\def\maxwidth{\ifdim\Gin@nat@width>\linewidth\linewidth\else\Gin@nat@width\fi}
\def\maxheight{\ifdim\Gin@nat@height>\textheight\textheight\else\Gin@nat@height\fi}
\makeatother
% Scale images if necessary, so that they will not overflow the page
% margins by default, and it is still possible to overwrite the defaults
% using explicit options in \includegraphics[width, height, ...]{}
\setkeys{Gin}{width=\maxwidth,height=\maxheight,keepaspectratio}
% Set default figure placement to htbp
\makeatletter
\def\fps@figure{htbp}
\makeatother

\IfFileExists{plex-otf.sty}{
  %% Full TeXlive
  % \usepackage[%
  %   % math,
  %   RM={Scale=0.94},SS={Scale=0.94},SScon={Scale=0.94},TT={Scale=MatchLowercase,FakeStretch=0.9},DefaultFeatures={Ligatures=Common}
  % ]{plex-otf}
}{
  %% TinyTeX
  \usepackage{libertine}
}

%%% Load theme
% https://deic.uab.cat/~iblanes/beamer_gallery/
\IfFileExists{beamerthemegotham.sty}{
  %% Full TeXlive
  \usetheme{gotham}
  \gothamset{
    numbering=totalpagenumber,
    parttocframe default=off,
    sectiontocframe default=off,
    subsectiontocframe default=off,
  }
}{
  %% TinyTeX
  \usetheme{Madrid}
}
\makeatletter
\@ifpackageloaded{caption}{}{\usepackage{caption}}
\AtBeginDocument{%
\ifdefined\contentsname
  \renewcommand*\contentsname{Содержание}
\else
  \newcommand\contentsname{Содержание}
\fi
\ifdefined\listfigurename
  \renewcommand*\listfigurename{Список иллюстраций}
\else
  \newcommand\listfigurename{Список иллюстраций}
\fi
\ifdefined\listtablename
  \renewcommand*\listtablename{Список таблиц}
\else
  \newcommand\listtablename{Список таблиц}
\fi
\ifdefined\figurename
  \renewcommand*\figurename{Рисунок}
\else
  \newcommand\figurename{Рисунок}
\fi
\ifdefined\tablename
  \renewcommand*\tablename{Таблица}
\else
  \newcommand\tablename{Таблица}
\fi
}
\@ifpackageloaded{float}{}{\usepackage{float}}
\floatstyle{ruled}
\@ifundefined{c@chapter}{\newfloat{codelisting}{h}{lop}}{\newfloat{codelisting}{h}{lop}[chapter]}
\floatname{codelisting}{Список}
\newcommand*\listoflistings{\listof{codelisting}{Листинги}}
\makeatother
\makeatletter
\makeatother
\makeatletter
\@ifpackageloaded{caption}{}{\usepackage{caption}}
\@ifpackageloaded{subcaption}{}{\usepackage{subcaption}}
\makeatother

\ifLuaTeX
\usepackage[bidi=basic]{babel}
\else
\usepackage[bidi=default]{babel}
\fi
\babelprovide[main,import]{russian}
\babelprovide[import]{english}
% get rid of language-specific shorthands (see #6817):
\let\LanguageShortHands\languageshorthands
\def\languageshorthands#1{}
\ifLuaTeX
  \usepackage{selnolig}  % disable illegal ligatures
\fi
\usepackage{csquotes}
\usepackage{bookmark}

\IfFileExists{xurl.sty}{\usepackage{xurl}}{} % add URL line breaks if available
\urlstyle{same} % disable monospaced font for URLs
\hypersetup{
  pdftitle={Лабораторная работа №4},
  pdfauthor={Назармамадов Умед Джамшедович},
  pdflang={ru-RU},
  hidelinks,
  pdfcreator={LaTeX via pandoc}}


\title{Лабораторная работа №4}
\subtitle{Агентная эпидемиологическая модель SIR на графе городов}
\author{Назармамадов Умед Джамшедович}
\date{2026-09-02}

\begin{document}
\frame{\titlepage}

\renewcommand*\contentsname{Содержание}
\begin{frame}[allowframebreaks]
  \frametitle{Содержание}
  \tableofcontents[hideallsubsections]
\end{frame}

\section{Информация}\label{ux438ux43dux444ux43eux440ux43cux430ux446ux438ux44f}

\begin{frame}{Докладчик}
\phantomsection\label{ux434ux43eux43aux43bux430ux434ux447ux438ux43a}
\begin{columns}[c]
\begin{column}{0.7\textwidth}
\begin{itemize}[<+->]
\tightlist
\item
  Назармамадов Умед Джамшедович
\item
  Российский университет дружбы народов им. П. Лумумбы
\item
  {[}1032225205@rudn.ru{]}
\item
  \url{https://umednazarmamadov.github.io/ru/}
\end{itemize}
\end{column}

\begin{column}{0.3\textwidth}
\begin{frame}{Цель работы}
\phantomsection\label{ux446ux435ux43bux44c-ux440ux430ux431ux43eux442ux44b}
\begin{itemize}[<+->]
\tightlist
\item
  Реализовать агентную модель SIR на графе городов
\item
  Учесть миграцию населения между городами
\item
  Исследовать влияние параметров на динамику эпидемии
\item
  Провести многокритериальную оптимизацию
\end{itemize}
\end{frame}

\begin{frame}{Агентная модель vs классическая SIR}
\phantomsection\label{ux430ux433ux435ux43dux442ux43dux430ux44f-ux43cux43eux434ux435ux43bux44c-vs-ux43aux43bux430ux441ux441ux438ux447ux435ux441ux43aux430ux44f-sir}
\begin{itemize}[<+->]
\tightlist
\item
  Классическая SIR --- система дифференциальных уравнений
\item
  Агентная модель --- каждый человек отдельный агент
\item
  Учитывается пространственная структура (граф городов)
\item
  Возможна миграция и локальные вспышки
\end{itemize}
\end{frame}

\begin{frame}{Реализация модели}
\phantomsection\label{ux440ux435ux430ux43bux438ux437ux430ux446ux438ux44f-ux43cux43eux434ux435ux43bux438}
\begin{itemize}[<+->]
\tightlist
\item
  Агент: статус (S/I/R), дни с момента заражения
\item
  Города --- вершины полного графа
\item
  Миграция по заданной матрице вероятностей
\item
  Заразность зависит от времени выявления инфекции
\end{itemize}
\end{frame}

\begin{frame}{Базовый эксперимент}
\phantomsection\label{ux431ux430ux437ux43eux432ux44bux439-ux44dux43aux441ux43fux435ux440ux438ux43cux435ux43dux442}
\includegraphics[width=0.6\textwidth,height=\textheight]{image/5.png}

Классическая картина развития эпидемии: рост, пик, спад числа
инфицированных.
\end{frame}

\begin{frame}{Сканирование параметра beta}
\phantomsection\label{ux441ux43aux430ux43dux438ux440ux43eux432ux430ux43dux438ux435-ux43fux430ux440ux430ux43cux435ux442ux440ux430-beta}
\includegraphics[width=0.6\textwidth,height=\textheight]{image/7.png}

Нелинейный рост пиковой заболеваемости и смертности с увеличением
заразности.
\end{frame}

\begin{frame}{Эффект миграции}
\phantomsection\label{ux44dux444ux444ux435ux43aux442-ux43cux438ux433ux440ux430ux446ux438ux438}
\includegraphics[width=0.6\textwidth,height=\textheight]{image/9.png}

Рост интенсивности миграции --- более ранний и высокий пик эпидемии.
\end{frame}

\begin{frame}{Многокритериальная оптимизация}
\phantomsection\label{ux43cux43dux43eux433ux43eux43aux440ux438ux442ux435ux440ux438ux430ux43bux44cux43dux430ux44f-ux43eux43fux442ux438ux43cux438ux437ux430ux446ux438ux44f}
\begin{itemize}[<+->]
\tightlist
\item
  Алгоритм Borg MOEA (BlackBoxOptim)
\item
  Критерии: пиковая заболеваемость и доля умерших
\item
  Параметры: beta, время выявления, смертность
\item
  Результат --- Парето-фронт компромиссных решений
\end{itemize}
\end{frame}

\begin{frame}{Результаты оптимизации}
\phantomsection\label{ux440ux435ux437ux443ux43bux44cux442ux430ux442ux44b-ux43eux43fux442ux438ux43cux438ux437ux430ux446ux438ux438}
\includegraphics[width=0.7\textwidth,height=\textheight]{image/11.png}

Раннее выявление снижает пик, но не всегда влияет на смертность
напрямую.
\end{frame}

\begin{frame}{Литературное программирование}
\phantomsection\label{ux43bux438ux442ux435ux440ux430ux442ux443ux440ux43dux43eux435-ux43fux440ux43eux433ux440ux430ux43cux43cux438ux440ux43eux432ux430ux43dux438ux435}
\begin{itemize}[<+->]
\tightlist
\item
  Все скрипты преобразованы через Literate.jl
\item
  Сгенерированы: чистый код, Quarto-документация, Jupyter Notebook
\item
  Notebooks выполнены и проверены
\end{itemize}
\end{frame}

\begin{frame}{Выводы}
\phantomsection\label{ux432ux44bux432ux43eux434ux44b}
\begin{itemize}[<+->]
\tightlist
\item
  Агентная модель SIR успешно реализована на графе городов
\item
  Показана нелинейная зависимость эпидемии от заразности
\item
  Миграция значимо ускоряет и усиливает эпидемию
\item
  Найдены компромиссные стратегии управления эпидемией
\end{itemize}
\end{frame}
\end{column}
\end{columns}
\end{frame}




\end{document}
